% Example "lotkavolterra": the problem, the solutions against the reference
% and the exact orbit, and the invariant tests. Every number comes from data/
% (parameters.tex, results.tex, orbit.dat, reference.dat,
% solution_<method>.dat, drift_<method>.dat, error_<method>.dat), written by
% examples/lotkavolterra/main.f90: run `fpm run --example lotkavolterra` from
% the project root, then build this file in examples/lotkavolterra/.
\documentclass[11pt,a4paper]{article}
\usepackage[margin=2.2cm]{geometry}
\usepackage{amsmath,booktabs}
\usepackage{pgfplots}
\usepgfplotslibrary{groupplots}
\pgfplotsset{compat=1.18}
\pagestyle{empty}

\begin{document}
\section*{Example: Lotka--Volterra predator--prey model}

\subsection*{1. Problem}
From the SciML problem library (\texttt{prob\_ode\_lotkavolterra}): prey $x$
grow and are eaten, predators $y$ die out and feed on the prey,
\begin{equation*}
  x' = a x - b x y, \qquad y' = -c y + d x y, \qquad x(0) = y(0) = 1 .
\end{equation*}
There is no closed form in time, but
\begin{equation*}
  V(x, y) = d x - c \ln x + b y - a \ln y
\end{equation*}
is constant along every solution ($\dot V = 0$), so the exact orbit is the
closed level curve $V = V_0$ around the equilibrium $(c/d, a/b)$, computed
here directly from $V$. In time, the solutions are compared with a
\texttt{dp54} run at a tiny step.
\begin{center}
\small
% Written by examples/lotkavolterra/main.f90: do not edit.
\begin{tabular}{ll}
\toprule
parameters & $(a, b, c, d) = (1.5, 1, 3, 1)$ \\
initial condition & $x(0) = 1$, $y(0) = 1$, $V_0 = 2$ \\
equilibrium & $(c/d, a/b) = (3, 1.5)$ \\
solution figure & $\Delta t = 0.1$, $0 \le t \le 20$ \\
reference curve & \texttt{dp54}, $\Delta t = 10^{-4}$ \\
convergence & $\max_t |V - V_0|$, $0 \le t \le 10$, $\Delta t = 0.1, \dots, 0.1/128$ \\
observed order & slope of the two finest levels with errors in $(10^{-12}, 1)$ \\
long time & $|V - V_0|$ for $0 \le t \le 1000$ at $\Delta t = 0.05$ \\
methods & \texttt{rk3ws}, \texttt{rk4}, \texttt{dp54} \\
interface & plain arrays, \texttt{rk\_array\_scheme\_t} \\
\bottomrule
\end{tabular}

\end{center}

\subsection*{2. Solutions}
\begin{center}
\begin{tikzpicture}
  \begin{groupplot}[
      group style={group size=2 by 1, horizontal sep=1.9cm}, % two panels side by side
      width=6.6cm, height=6.2cm,                           % size of each panel
      grid=major, grid style={black!12},                   % light grid
    ]
    % ---- (a) populations in time ----
    \nextgroupplot[
      title={(a) $x(t)$, $y(t)$ at $\Delta t = 0.1$},      % panel title
      xlabel={$t$}, xmin=0, xmax=20, ymin=0, ymax=10,      % time, populations
      legend to name=lvlegend, legend columns=5,           % legend below both panels
      legend style={font=\small, /tikz/every even column/.append style={column sep=0.3cm}},
    ]
    \addplot[black, thick] table[x=t, y=x] {data/reference.dat};                  % reference prey
    \addlegendentry{exact orbit / reference}
    \addplot[black, thick, dashed, forget plot] table[x=t, y=y] {data/reference.dat}; % reference predators
    \addplot[blue!80!black, thin] table[x=t, y=x] {data/solution_rk4.dat};        % RK4 prey
    \addlegendentry{\texttt{rk4}}
    \addplot[blue!80!black, thin, dashed, forget plot] table[x=t, y=y] {data/solution_rk4.dat}; % RK4 predators
    % ---- (b) phase plane ----
    \nextgroupplot[
      title={(b) phase plane, $\Delta t = 0.1$},           % panel title
      xlabel={prey $x$}, ylabel={predators $y$},           % orbit in the (x, y) plane
      xmin=0, xmax=9, ymin=0, ymax=6,                      % around the orbit
    ]
    \addplot[blue!80!black, thin] table[x=x, y=y] {data/solution_rk4.dat};    % RK4: on the orbit
    \addplot[black, thick] table[x=x, y=y] {data/orbit.dat};                  % exact orbit V = V0
    \addplot[only marks, mark=*, mark size=1.5pt] coordinates {(3, 1.5)};     % equilibrium (c/d, a/b)
  \end{groupplot}
  \node[below=0.2cm] at (current bounding box.south)        % legend centred under both panels
    {\pgfplotslegendfromname{lvlegend}};
\end{tikzpicture}
\end{center}
(a) Prey (solid) and predators (dashed) oscillate with a period of about 5.
At $\Delta t = 0.1$ \texttt{rk4} follows the reference. (b) In the phase
plane the exact solution runs around the closed curve $V = V_0$ (the dot is
the equilibrium), and \texttt{rk4} stays on it at this scale.

\newpage
\subsection*{3. Tests: the invariant $V$}

\paragraph{Prediction.} Runge--Kutta methods do not keep nonlinear invariants
such as $V$; the drift $|V - V_0|$ is a smooth function of the global error,
so (i) it converges at the method's (general) order as $\Delta t$ shrinks,
and (ii) at a fixed step it has two parts: an oscillation within each cycle,
bounded, and a secular part that grows linearly in time as the orbit moves
from one level curve to the next.

\paragraph{Result.}
\begin{center}
\small
% Written by examples/lotkavolterra/main.f90: do not edit.
\begin{tabular}{llcccc}
\toprule
method & name & order & observed & $|V - V_0|$ at $t = 1000$ & positive until \\
\midrule
RK3 of Wicker and Skamarock & \texttt{rk3ws} & 2 &  2.00 & $2.13\times10^{-2}$ & end \\
Classic RK4 & \texttt{rk4} & 4 &  3.99 & $2.37\times10^{-3}$ & end \\
Dormand-Prince 5(4) & \texttt{dp54} & 5 &  4.77 & $9.62\times10^{-6}$ & end \\
\bottomrule
\end{tabular}

\end{center}

\begin{center}
\begin{tikzpicture}
  \begin{groupplot}[
      group style={group size=2 by 1, horizontal sep=1.9cm}, % two panels side by side
      width=6.6cm, height=6.2cm,                           % size of each panel
      grid=major, grid style={black!12},                   % light grid
      every axis plot/.append style={thick},               % same weight for all curves
    ]
    % ---- (i) convergence ----
    \nextgroupplot[
      xmode=log, ymode=log,                                % both axes logarithmic
      title={(i) $\max_t |V - V_0|$, $0 \le t \le 10$},    % panel title
      xlabel={$\Delta t$},                                 % step
      legend to name=lvtestlegend, legend columns=7,       % legend below both panels
      legend style={font=\small, /tikz/every even column/.append style={column sep=0.2cm}},
    ]
    \addplot[green!50!black, dashed, mark=triangle, mark size=1.5pt] table[x=dt, y=error] {data/error_rk3ws.dat}; % order 2 here
    \addlegendentry{\texttt{rk3ws}}
    \addplot[blue!80!black, mark=diamond*, mark size=1.5pt]   table[x=dt, y=error] {data/error_rk4.dat};      % order 4
    \addlegendentry{\texttt{rk4}}
    \addplot[red!75!black, mark=pentagon*, mark size=1.5pt]   table[x=dt, y=error] {data/error_dp54.dat};     % order 5
    \addlegendentry{\texttt{dp54}}
    % ---- (ii) long time ----
    \nextgroupplot[
      xmode=log, ymode=log,                                % both axes logarithmic
      title={(ii) $\max_{s \le t}|V - V_0|$, $\Delta t = 0.05$}, % panel title: running maximum
      xlabel={$t$}, ymin=1e-12,                            % time up to t_long
    ]
    \addplot[green!50!black, dashed] table[x=t, y=drift] {data/drift_rk3ws.dat};   % order 2 here
    \addplot[blue!80!black]         table[x=t, y=drift] {data/drift_rk4.dat};      % order 4
    \addplot[red!75!black]          table[x=t, y=drift] {data/drift_dp54.dat};     % order 5
  \end{groupplot}
  \node[below=0.2cm] at (current bounding box.south)        % legend centred under both panels
    {\pgfplotslegendfromname{lvtestlegend}};
\end{tikzpicture}
\end{center}
(i) The drift converges at each method's general order; \texttt{rk3ws} is of
order 2 on this nonlinear problem, and \texttt{dp54} is still approaching
order 5 from above (as in the nonlinear chemistry example). (ii) The running
maximum of the drift: for \texttt{rk4} and \texttt{dp54} it grows tenfold
per decade of $t$ from $t \approx 10$ on (the steps below are single cycles),
the linear secular part; for \texttt{rk3ws}, of lower order, the oscillation
within a cycle dominates first and the linear growth takes over later. No
explicit Runge--Kutta method stays on the orbit forever; all three stay
bounded over the whole run.
\end{document}
